% Use a modern class instead of an old article class
\documentclass{scrartcl}

% The class "scrartcl" provides the “article”-like element of the koma-script collection. The document layout of the class is less 'strident' than that of article, and it offers much more flexibility than article via other elements of the koma-script collection.

% Information about other document classes can be found herehttps://en.wikibooks.org/wiki/LaTeX/Document_Structure#Document_classes


% %%%%%%%%%%%%%%%%%%%%%%%%%%% CUSTOM PACKAGES %%%%%%%%%%%%%%%%%%%%%%%%%%% %
\usepackage[none]{hyphenat}
% \usepackage{geometry}
\usepackage{amssymb}
\usepackage{amsthm}
\usepackage{empheq}
\usepackage{mdframed}
\usepackage{lipsum}
\usepackage{color}
\usepackage{bm}
\usepackage{array}
\usepackage{enumitem}
\usepackage{subfigure}
\usepackage[most]{tcolorbox}

\usepackage[a4paper, left=2cm, right=2cm]{geometry}

% Long table
\usepackage{xltabular}

%maths
\usepackage{amsfonts}

\usepackage[edges]{forest}
\usepackage[T1]{fontenc}

%tikzpicture
\usepackage{tikz}
\usepackage{scalerel}
\usepackage{tkz-euclide}
\usetikzlibrary{trees}
\usetikzlibrary{calc}
\usetikzlibrary{patterns,arrows.meta}
\usetikzlibrary{shadows}
\usetikzlibrary{external}
\usepackage{tikz-cd}


%pgfplots
\usepackage{pgfplots}  % For plots and graphs
\usepgfplotslibrary{statistics}
\usepgfplotslibrary{fillbetween}

% \usepackage{cellspace} % Adds extra space inside table cells
% \setlength\cellspacetoplimit{12px}    % Top padding
% \setlength\cellspacebottomlimit{12px} % Bottom padding

\renewcommand{\arraystretch}{1.5} % Increase row height (default = 1.0)
\usepackage{tabularx}
% Define a new column type: vertically centered and flexible width
\newcolumntype{M}[1]{>{\centering\arraybackslash}m{#1}}
\newcolumntype{Y}{>{\centering\arraybackslash}X} % Like M but auto width
% %%%%%%%%%%%%%%%%%%%%%%%%%%% CUSTOM PACKAGES %%%%%%%%%%%%%%%%%%%%%%%%%%% %
% %%%%%%%%%%%%%%%%%%%%%%%%%%% CUSTOM COMMANDS %%%%%%%%%%%%%%%%%%%%%%%%%%% %
% \usepackage{tocbasic}
% \DeclareNewSectionCommand[
%   level=4,
%   indent=0pt,
%   beforeskip=1.5ex plus .2ex minus .1ex,
%   afterskip=.8ex plus .2ex,
%   style=section,
%   font=\normalfont\normalsize\bfseries
% ]{subsubsubsection}

% % TOC depth and numbering depth
% \setcounter{tocdepth}{4}
% \setcounter{secnumdepth}{4}

% % Adjust indentation in TOC
% \DeclareTOCStyleEntry[
%   indent=1.5em,
%   numwidth=3em
% ]{tocline}{subsubsubsection}

% Extend sectioning to level 4 (subsubsubsection)
\makeatletter
\newcounter{subsubsubsection}[subsubsection]
\renewcommand\thesubsubsubsection{\thesubsubsection.\arabic{subsubsubsection}}

\newcommand\subsubsubsection{\@startsection{subsubsubsection}{4}{\z@}%
  {-3.25ex \@plus -1ex \@minus -.2ex}%
  {1.5ex \@plus .2ex}%
  {\normalfont\normalsize\bfseries}}

\newcommand\l@subsubsubsection{\@dottedtocline{4}{5em}{4em}}
\makeatother

% Add subsubsubsection to TOC and numbering
\setcounter{secnumdepth}{4}
\setcounter{tocdepth}{4}

% %%%%%%%%%%%%%%%%%%%%%%%%%%% CUSTOM COMMANDS %%%%%%%%%%%%%%%%%%%%%%%%%%% %
% --------------------------- DEFAULT PACKAGES SETUP --------------------------- %
\usepackage[utf8]{inputenc} % Specify the input encoding of the source file as UTF8 
\usepackage[T1]{fontenc} % The T1 font encoding is an 8-bit encoding and fully supports words containing accented characters
\usepackage{lmodern} % Load Latin modern fonts in outline format
\usepackage{inconsolata} % Load better typewriter font
\usepackage[english]{babel} % For better hyphenation patterns
\usepackage{microtype} % For improved micro-typography
\frenchspacing % Switch off extra space after a sentence
\usepackage{graphicx} % To be able to iclude pictures
\usepackage{mathtools,amsmath} % for advanced math formulas
\usepackage{minted} % To highlight source code
% Default options for all minted environments
\setminted{
    style=default, % for minted envronment
    autogobble, % automatically remove all common leading whitespace
    xleftmargin=10pt, % indentation to add (on the left) before the listing
    numbersep=6pt, % gap between numbers and start of line
    linenos, % enables line numbers
    mathescape % enables the usual math mode inside comments
}
\setmintedinline[c]{style=default} % for minted inlines
\usepackage{fancyvrb} % Flexible handling of verbatim text
\usepackage{booktabs} % For well-spaced lines and guidelines in tables
\usepackage{algorithm} % To typeset algorithms or pseudocode in LaTeX 
\usepackage{algpseudocode}
\pgfplotsset{width=10cm,compat=1.9}
\renewcommand{\algorithmiccomment}[1]{\hspace{1cm}// #1}

\usepackage{soul} % Change text highlighting
\usepackage{pgffor} % Needed for for-loops

% --------------------------- DEFAULT PACKAGES SETUP --------------------------- %
% --------------------------- KOMA OPTIONS (DOCUMENT SETTINGS) ----------------------- %

\KOMAoptions{
    parskip=half,  % full, off
    fontsize=12pt, % base font size (10pt default)
    toc=sectionentrydotfill, % Fill toc with dots for sections
    % headings=big,% small/normal/big headings (normal is default), 
    % paper=a5,    % paper format (a4 default) 
    pagesize=auto  % Use paper format for PDF too
}

% --------------------------- KOMA OPTIONS (DOCUMENT SETTINGS) ----------------------- %
%%%%%%%%%%%%%%%%%%%%%%%%%% Define some useful colors %%%%%%%%%%%%%%%%%%%%%%%%%%
\definecolor{ocre}{RGB}{243,102,25}
\definecolor{mygray}{RGB}{243,243,244}
\definecolor{deepGreen}{RGB}{26,111,0}
\definecolor{shallowGreen}{RGB}{235,255,255}
\definecolor{deepBlue}{RGB}{61,124,222}
\definecolor{shallowBlue}{RGB}{235,249,255}
%%%%%%%%%%%%%%%%%%%%%%%%%%%%%%%%%%%%%%%%%%%%%%%%%%%%%%%%%%%%%%%%%%%%%%%%%%%%%%%

%%%%%%%%%%%%%%%%%%%%%%%%%% Define an orangebox command %%%%%%%%%%%%%%%%%%%%%%%%
\newcommand\orangebox[1]{\fcolorbox{ocre}{mygray}{\hspace{1em}#1\hspace{1em}}}
%%%%%%%%%%%%%%%%%%%%%%%%%%%%%%%%%%%%%%%%%%%%%%%%%%%%%%%%%%%%%%%%%%%%%%%%%%%%%%%

%%%%%%%%%%%%%%%%%%%%%%%%%%%% English Environments %%%%%%%%%%%%%%%%%%%%%%%%%%%%%
\newtheoremstyle{mytheoremstyle}{3pt}{3pt}{\normalfont}{0cm}{\rmfamily\bfseries}{}{1em}{{\color{black}\thmname{#1}~\thmnumber{#2}}\thmnote{\,--\,#3}}
\newtheoremstyle{myproblemstyle}{3pt}{3pt}{\normalfont}{0cm}{\rmfamily\bfseries}{}{1em}{{\color{black}\thmname{#1}~\thmnumber{#2}}\thmnote{\,--\,#3}}
\theoremstyle{mytheoremstyle}
\newmdtheoremenv[linewidth=1pt,backgroundcolor=shallowGreen,linecolor=deepGreen,leftmargin=0pt,innerleftmargin=20pt,innerrightmargin=20pt,]{theorem}{Theorem}[section]
\theoremstyle{mytheoremstyle}
\newmdtheoremenv[linewidth=1pt,backgroundcolor=shallowBlue,linecolor=deepBlue,leftmargin=0pt,innerleftmargin=20pt,innerrightmargin=20pt,]{definition}{Definition}[section]
\theoremstyle{myproblemstyle}
\newmdtheoremenv[linecolor=black,leftmargin=0pt,innerleftmargin=10pt,innerrightmargin=10pt,]{problem}{Problem}[section]
%%%%%%%%%%%%%%%%%%%%%%%%%%%%%%%%%%%%%%%%%%%%%%%%%%%%%%%%%%%%%%%%%%%%%%%%%%%%%%%

%%%%%%%%%%%%%%%%%%%%%%%%%%%%%%% Plotting Settings %%%%%%%%%%%%%%%%%%%%%%%%%%%%%
\usepgfplotslibrary{colorbrewer}
\pgfplotsset{width=8cm,compat=1.9}
%%%%%%%%%%%%%%%%%%%%%%%%%%%%%%%%%%%%%%%%%%%%%%%%%%%%%%%%%%%%%%%%%%%%%%%%%%%%%%%
% %%%%%%%%%%%%%%%%%%%%%%%%%%%%%%%%%%%%%% REPORT DETAILS %%%%%%%%%%%%%%%%%%%%%%%%%%%%%% %

\newcommand{\courseyear}{2026}
\newcommand{\coursename}{TFSS25 - Reinforcement Learning}
\newcommand{\coursecredits}{7.5 credits}
\newcommand{\fullcoursename}{\coursename, \coursecredits}
\newcommand{\termname}{Autumn\courseyear}
% \newcommand{\groupnumber}{10} % Optional
\newcommand{\reportname}{Reinforcement Learning}
\newcommand{\assignmentname}{Reinforcement Learning}

% %%%%%%%%%%%%%%%%%%%%%%%%%%%%%%%%%% REPORT DETAILS %%%%%%%%%%%%%%%%%%%%%%%%%%%%%%%%%% %
% -------------------------- SETTING UP PROPERTIES OF PDF FILE ----------------------- %

\usepackage{hyperref}
\usepackage[all]{hypcap}
\hypersetup{
    colorlinks=true,
    allcolors=blue,
    pdftitle={\reportname},
    pdfauthor={Martin Nilsson},
    pdfsubject={\coursename, \termname}
	% pdfpagemode=FullScreen,
	% pdfpagemode= UseOutlines % the default if present
}

% ---------------------------- SETTING UP PROPERTIES OF PDF FILE ------------------------ %
% ----------------------------- HEADER AND FOOTER SETTINGS ------------------------------ %

\usepackage{scrlayer-scrpage}
\usepackage{zref-totpages} % To get the total number of pages
\pagestyle{scrheadings}
\KOMAoptions{% 
    headsepline,% line below the header
    plainheadsepline,% also on scrplain
}
\clearscrheadfoot % Erase all current configurations
% inner part of the header [titel_page]{other_pages}
\ihead[\coursename, \courseyear]{\coursename, \courseyear}
% outer part of the header [titel_page]{other_pages}
\ohead{\assignmentname}
% center part of the footer [titel_page]{other_pages}
\cfoot[\textup{Page \pagemark\ of \ztotpages}]{\textup{Page \pagemark\ of \ztotpages}}

% ----------------------------- HEADER AND FOOTER SETTINGS ------------------------------ %
% %%%%%%%%%%%%%%%%%%%%%%%%%%%%%%%%%% TITLE PAGE CONFIG %%%%%%%%%%%%%%%%%%%%%%%%%%%%%%%%%% %

\title{\reportname}
\subtitle{Assignment 2 - Lunar Lander}
\author{
    \textbf{Authors:} \\
    \textbf{Martin Nilsson}, \textit{School of Engineering}, Jönköping University
    \and
    \textbf{Besan Ewir}, \textit{School of Engineering}, Jönköping University
    \and
    \textbf{Vladimir Giustacchini}, \textit{School of Engineering}, Jönköping University
    \and
    \textbf{Jorge Cuello}, \textit{School of Engineering}, Jönköping University
}
\date{\today}

% %%%%%%%%%%%%%%%%%%%%%%%%%%%%%%%%%% TITLE PAGE CONFIG %%%%%%%%%%%%%%%%%%%%%%%%%%%%%%%%%% %


\begin{document} % <<<<<<<<<<<<<<<<<<<<<<<<<<<<<<<<<<<<<<<<<<<<<<
\raggedright

\maketitle

\tableofcontents

\newpage

\section{Theoretical Part}

\subsection{Question 1}

\subsubsection*{\textcolor{blue}{Q1: ``Explain the concept of the Q-learning update rule. How does it relate to the Bellman optimality equation? What are the potential challenges in using this update rule? ``}}

\subsection{Question 2}

\subsubsection*{\textcolor{blue}{Q2: ``Discuss the role of experience replay in DQN. How does it help to address the correlation problem in reinforcement learning? ``}}
Experience Replay is one of the critical elements of  Deep Q-Networks (DQN),  where the agent saves its experiences  $(s, a, r, s')$ in the memory, which can be used during learning rather than learning from the latest experience only. In the training process, samples are randomly selected from the experience replay buffer.

The main benefit of using experience replay is that it helps reduce the correlation problem,  since in reinforcement learning, the experiences that come one after another are highly correlated. Random sampling mixes experiences from different time steps, which reduces the correlation between successive updates and makes learning more stable.

Experience Replay also allows the same experience to be used multiple times while the agent still learning, making learning more efficient and reducing the variance of the updates.

\subsection{Question 3}

\subsubsection*{\textcolor{blue}{Q3: ``Explain the concept of the target network in DQN. How does it help to stabilize the learning process? ``}}

\subsubsection*{Target Network:} The target network is a copy of the Q-network that is used to calculate the target Q-values during training. The target network is updated less frequently than the Q-network, which helps to stabilize the learning process by reducing the correlation between the Q-values and the target values.

Keep in mind that the \textit{target network} is \textbf{NOT} necessarily a different architecture or a separate trained network. In most cases it is simply a periodically updated copy of the online network.

In a neural network with $\theta$ parameters, we have $Q(s, a; \theta)$ and the target becomes $$y = r + \gamma \max_{a'} Q(s', a'; \theta)$$

But since we use the same network to calculate the target, it can lead to instability and divergence in the learning process because the target would move with the parameters. To address this issue, we introduce a \textit{target network} with parameters $\theta^-$, which is a copy of the Q-network that is updated less frequently (e.g., every $C$ steps). The target becomes $$y = r + \gamma \max_{a'} Q(s', a'; \theta^-)$$

\begin{itemize}
    \item $Q(s, a; \theta)$ (\textbf{Online network} $\rightarrow$ the network that is being trained)
    \item $Q(s', a'; \theta^-)$ (\textbf{Target network} $\rightarrow$ the network that is used to calculate the target Q-values)
\end{itemize}


A simplified DQN training loop with a target network can be described as follows:
\begin{tcolorbox}[
        colback=blue!5!white,
        colframe=blue!60!black,
        title=\textbf{DQN Target-Network Pseudocode},
        sharp corners,
        boxrule=0.8pt
    ]
    \textbf{Initialize} online network $\theta$\\
    \textbf{Initialize} target network $\theta^- = \theta$\\[2pt]

    \textbf{Repeat:}
    \begin{enumerate}[leftmargin=*, itemsep=2pt]
        \item Take an action using the online network.
        \item Store transition $(s, a, r, s')$ in the Experience/Replay buffer.
        \item Sample a mini-batch from the replay buffer.
        \item Calculate targets using the target network:
              \[
                  y = r + \gamma \max_{a'} Q(s', a'; \theta^-).
              \]
        \item Update the online network $\theta$ to minimize:
              \[
                  \bigl(y - Q(s,a;\theta)\bigr)^2.
              \]
        \item Every $C$ steps, update the target network:
              \[
                  \theta^- \leftarrow \theta.
              \]
    \end{enumerate}
\end{tcolorbox}

The main takeaway is that the target network helps to stabilize the learning process by providing a more stable target for the Q-value updates, which reduces the risk of divergence and improves convergence.

% =========== REFERENCES =========== 

\bibliographystyle{ieeetr}
\bibliography{rl}

\end{document}